\documentclass[11pt]{article}

\usepackage[utf8]{inputenc}
\usepackage[T1]{fontenc}
\usepackage{geometry}
\usepackage{amsmath}
\usepackage{amssymb}
\usepackage{booktabs}
\usepackage{hyperref}
\usepackage{xurl}
\usepackage{xcolor}
\usepackage{listings}
\usepackage{enumitem}
\usepackage{parskip}
\usepackage{graphicx}
\graphicspath{{../images/}}

\geometry{margin=1in}

\lstset{
  language=Java,
  basicstyle=\ttfamily\small,
  keywordstyle=\color{blue},
  commentstyle=\color{gray},
  stringstyle=\color{teal},
  showstringspaces=false,
  columns=fullflexible,
  frame=single,
  framerule=0.2pt,
  breaklines=true
}

\setlist[itemize]{topsep=0.3em,itemsep=0.2em}
\setlist[enumerate]{topsep=0.3em,itemsep=0.2em}

% Simple per-section source line.
\newcommand{\notesources}[1]{\par\noindent{\small\color{gray}\textit{Sources:} \sloppy #1\par}}

% Unit-wise source strings for this subject (used by slides/notes).
% Path is relative to the lecture `latex/` folder (the compilation CWD).
% OOPS with Java (CSEG1044) — unit-wise sources
%
% These are short, student-facing reference strings intended to be used with:
%   - \slidesources{...}  (in beamer slides)
%   - \notesources{...}   (in lecture notes)
%
% Style inspiration: other subjects in My-Subjects/ use semicolon-separated sources
% and include an "accessed" date for web documentation.

\newcommand{\OOPSUnitISources}{Schildt, \textit{Java: A Beginner's Guide} (10e), McGraw-Hill (Java fundamentals + OOP basics).; Oracle Java Tutorials: Getting Started + Language Basics + Classes and Objects (accessed \today).; Java SE API docs (e.g., \texttt{String}, \texttt{Scanner}) (accessed \today).}

\newcommand{\OOPSUnitIISources}{Schildt, \textit{Java: The Complete Reference} (12e), McGraw-Hill (classes, inheritance, packages, interfaces).; Bloch, \textit{Effective Java} (3e), Addison-Wesley, 2018 (design choices; interfaces vs abstract classes).; Oracle Java Tutorials: Classes and Objects + Interfaces + Packages (accessed \today).; Java SE API docs (e.g., \texttt{Comparable}, \texttt{Runnable}) (accessed \today).}

\newcommand{\OOPSUnitIIISources}{Schildt, \textit{Java: The Complete Reference} (12e) (nested/inner classes, exceptions, threads, I/O).; Oracle Java Tutorials: Nested Classes + Exceptions + Concurrency + Basic I/O (accessed \today).; Java SE API docs (e.g., \texttt{Thread}, \texttt{AutoCloseable}) (accessed \today).}

\newcommand{\OOPSUnitIVSources}{Schildt, \textit{Java: The Complete Reference} (12e) (generics, lambdas, Swing, JDBC).; Bloch, \textit{Effective Java} (3e) (generics + lambdas).; Oracle Java Tutorials: Generics + Lambda Expressions + Swing + JDBC Basics (accessed \today).; Java SE API docs (e.g., \texttt{java.util.function}, \texttt{javax.swing}, \texttt{java.sql}) (accessed \today).}

\newcommand{\OOPSUnitVSources}{Schildt, \textit{Java: The Complete Reference} (12e) (Collections Framework + wrapper classes).; Oracle Java docs: Collections Framework overview (accessed \today).; Java SE API docs (\texttt{java.util} collections; wrapper classes in \texttt{java.lang}) (accessed \today).}

\newcommand{\OOPSUnitVISources}{Git SCM: \textit{Pro Git} (2e) (version control workflow), accessed \today.; GitHub Docs (repositories, issues, pull requests), accessed \today.; Oracle Java Tutorials: Swing + JDBC (accessed \today).; JUnit 5 User Guide (accessed \today).; Apache Maven Project: Getting Started (accessed \today).; Gradle User Manual: Getting Started (accessed \today).; UML class diagrams: UML 2.x overview/spec (accessed \today).}

\newcommand{\OOPSMiscWhyInterfacesSources}{Schildt, \textit{Java: The Complete Reference} (12e) (interfaces).; Bloch, \textit{Effective Java} (3e) (prefer interfaces where appropriate).; Oracle Java Tutorials: Interfaces (accessed \today).; Java SE API docs (e.g., \texttt{Comparable}, \texttt{Runnable}) (accessed \today).}



\title{Object Oriented Programming (CSEG1044)\\Unit 04 -- Lecture 01 Notes\\Generics Fundamentals (Generic Classes)}
\author{Tofik Ali}
\date{\today}

\begin{document}
\maketitle
\tableofcontents

\section{Lecture Overview}
Generics help you write reusable code without losing type safety.
The main idea is simple:
\begin{quote}
Write a class/method once, and let the compiler enforce the types for each use.
\end{quote}
In this lecture we focus on \textbf{generic classes} like \texttt{Box<T>} and learn why raw types are dangerous.
\notesources{\OOPSUnitIVSources}

\section{Syllabus Alignment (Unit 04)}
\textbf{Unit 04:} Generics fundamentals; generic classes

\section{Learning Outcomes}
After this lecture, you should be able to:
\begin{itemize}
  \item Explain the motivation of generics (type safety + reuse) with a small example.
  \item Create and use a generic class such as \texttt{Box<T>}.
  \item Identify raw-type usage and explain why it can cause runtime errors.
\end{itemize}

\section{Key Terms (Quick)}
\begin{itemize}
  \item \textbf{Generic}: code with a type parameter (e.g., \texttt{<T>}).
  \item \textbf{Type parameter}: placeholder type name like \texttt{T}, \texttt{K}, \texttt{V}.
  \item \textbf{Parameterized type}: using a generic with a specific type, e.g., \texttt{Box<String>}.
  \item \textbf{Raw type}: using a generic without type arguments, e.g., \texttt{Box} (unsafe).
  \item \textbf{Type erasure (idea)}: generics are mainly compile-time; runtime types are erased (advanced detail).
\end{itemize}

\section{Detailed Notes}
\subsection{Why generics? (the problem without generics)}
Before generics (or when you avoid them), you often use \texttt{Object}:
\begin{itemize}
  \item You can store anything in \texttt{Object}.
  \item But you lose type information and must cast back.
  \item Wrong casts fail at runtime with \texttt{ClassCastException}.
\end{itemize}

\textbf{What generics give you:}
\begin{itemize}
  \item compile-time checking (errors caught earlier),
  \item no (or less) casting,
  \item clearer code: \texttt{List<String>} tells you what the list contains.
\end{itemize}

\subsection{Generic classes: \texttt{Box<T>}}
A generic class has a type parameter:
\begin{itemize}
  \item \texttt{class Box<T> \{ ... \}}
\end{itemize}
\texttt{T} can represent any reference type (like \texttt{String}, \texttt{Student}, \texttt{Integer}).

\paragraph{Naming convention.}
\begin{itemize}
  \item \texttt{T} = Type,\;\; \texttt{E} = Element,\;\; \texttt{K/V} = Key/Value.
\end{itemize}

\subsection{Avoiding raw types (common mistakes)}
Raw type means you wrote \texttt{Box} instead of \texttt{Box<T>}. It compiles with warnings but is unsafe.

\textbf{Safer alternative:}
\begin{itemize}
  \item If you truly do not know the type, use \texttt{Box<?>} rather than raw \texttt{Box}.
\end{itemize}

\textbf{Common mistakes:}
\begin{itemize}
  \item Using raw collections: \texttt{List list = new ArrayList();}
  \item Mixing types inside one list, then casting incorrectly.
  \item Confusing generics with primitives: use \texttt{Integer} not \texttt{int} in \texttt{List<int>} (not allowed).
  \item Thinking generics exist at runtime the same way: due to type erasure, some runtime checks are not possible (advanced).
\end{itemize}
\notesources{\OOPSUnitIVSources}

\section{Worked Example: \texttt{Box<T>} + raw type danger}
\subsection*{Generic class}
\begin{lstlisting}
class Box<T> {
    private final T value;
    Box(T value) { this.value = value; }
    T get() { return value; }
}
\end{lstlisting}

\subsection*{Safe usage (compile-time checks)}
\begin{lstlisting}
public class DemoSafe {
    public static void main(String[] args) {
        Box<String> b1 = new Box<>("UPES");
        String s = b1.get(); // no cast

        Box<Integer> b2 = new Box<>(10);
        int x = b2.get(); // auto-unboxing
        System.out.println(s + " " + x);
    }
}
\end{lstlisting}

\subsection*{Unsafe raw type (runtime failure risk)}
\begin{lstlisting}
public class DemoRaw {
    public static void main(String[] args) {
        Box raw = new Box(123);  // raw type: warning
        String s = (String) raw.get(); // ClassCastException at runtime
        System.out.println(s);
    }
}
\end{lstlisting}

\section{In-Class Exercises (with brief solutions)}
\begin{enumerate}
  \item Why is \texttt{Box} (raw type) unsafe?\par
  \textbf{Answer:} it loses type checking and can lead to wrong casts at runtime.

  \item What does \texttt{Box<String>} mean?\par
  \textbf{Answer:} a box that can hold only \texttt{String} values (type-safe).

  \item Can we write \texttt{List<int>}?\par
  \textbf{Answer:} no, generics use reference types; use \texttt{List<Integer>}.
\end{enumerate}

\section{Self-Study Practice (no solutions)}
\begin{enumerate}
  \item Create a generic class \texttt{Pair<K,V>} with getters.
  \item Write a program that stores \texttt{Student} objects in \texttt{List<Student>} and prints them.
  \item Find one example of raw type in your old code and refactor it to a parameterized type.
\end{enumerate}

\section{Checkpoint Questions (with sample answers)}
\textbf{Q1.} What problem do generics solve?\par
\textbf{Sample answer:} Generics provide compile-time type safety for reusable code. They reduce casts and prevent runtime \texttt{ClassCastException} caused by mixing types.

\vspace{0.6em}
\textbf{Q2.} Why should we avoid raw types?\par
\textbf{Sample answer:} Raw types disable generic type checking and allow unsafe operations, which can cause runtime failures and reduce code readability.

\end{document}
